% !TEX program = xelatex
% =====================================================================
% LLM Theory Seminar -- English Study Notes
% Compile with:
%   xelatex FILENAME.tex
%   xelatex FILENAME.tex
% or: latexmk -xelatex FILENAME.tex
% =====================================================================
\documentclass[11pt,a4paper,oneside,openany]{memoir}
\usepackage{amsmath,amssymb,amsthm,mathtools,bm}
\usepackage{booktabs,longtable,array,tabularx,makecell}
\usepackage{enumitem}
\usepackage[most]{tcolorbox}
\usepackage[dvipsnames]{xcolor}
\usepackage[unicode,bookmarks=false]{hyperref}
\usepackage{cleveref}
\usepackage{needspace}
\setlrmarginsandblock{27mm}{27mm}{*}
\setulmarginsandblock{24mm}{28mm}{*}
\checkandfixthelayout
\setlength{\parindent}{0pt}
\setlength{\parskip}{0.48em}
\setlength{\emergencystretch}{3em}
\hbadness=10000
\setlength{\headheight}{15pt}
\linespread{1.055}\selectfont
\setlist[itemize]{topsep=0.28em,itemsep=0.12em,parsep=0pt,leftmargin=1.55em}
\setlist[enumerate]{topsep=0.28em,itemsep=0.16em,parsep=0pt,leftmargin=1.75em}
\definecolor{NoteBlue}{HTML}{294E6B}
\definecolor{NoteBlueLight}{HTML}{F2F6F9}
\definecolor{NoteGray}{HTML}{5A6268}
\definecolor{NoteGrayLight}{HTML}{F6F6F4}
\definecolor{NoteWarm}{HTML}{7A6545}
\definecolor{NoteWarmLight}{HTML}{FAF7F0}
\hypersetup{colorlinks=true,linkcolor=NoteBlue,citecolor=NoteBlue,urlcolor=NoteBlue}
\makepagestyle{seminar}
\makeoddhead{seminar}{\footnotesize LLM Theory Seminar}{}{\footnotesize In-Context Learning Theory}
\makeoddfoot{seminar}{}{\footnotesize\thepage}{}
\makeheadrule{seminar}{\textwidth}{0.35pt}
\pagestyle{seminar}
\aliaspagestyle{chapter}{seminar}
\makechapterstyle{seminarnotes}{
  \setlength{\beforechapskip}{8pt}
  \setlength{\midchapskip}{8pt}
  \setlength{\afterchapskip}{22pt}
  \renewcommand{\chapterheadstart}{\vspace*{\beforechapskip}}
  \renewcommand{\chapnamefont}{\normalfont\small\bfseries}
  \renewcommand{\chapnumfont}{\normalfont\bfseries\fontsize{34}{38}\selectfont\color{NoteBlue}}
  \renewcommand{\chaptitlefont}{\normalfont\huge\bfseries\color{black}}
  \renewcommand{\printchaptername}{}
  \renewcommand{\chapternamenum}{}
  \renewcommand{\printchapternum}{\chapnumfont\thechapter}
  \renewcommand{\afterchapternum}{\par\nobreak\color{black}\vskip\midchapskip}
  \renewcommand{\printchaptertitle}[1]{\chaptitlefont ##1}
  \renewcommand{\afterchaptertitle}{\par\nobreak\vskip\afterchapskip}
}
\chapterstyle{seminarnotes}
\setsecheadstyle{\Large\bfseries}
\setsubsecheadstyle{\large\bfseries}
\setsubsubsecheadstyle{\normalsize\bfseries}
\setbeforesecskip{-1.3\baselineskip plus -0.2\baselineskip minus -0.1\baselineskip}
\setaftersecskip{0.55\baselineskip}
\setbeforesubsecskip{-0.9\baselineskip plus -0.15\baselineskip minus -0.1\baselineskip}
\setaftersubsecskip{0.35\baselineskip}
\newcommand{\R}{\mathbb{R}}
\newcommand{\E}{\mathbb{E}}
\newcommand{\norm}[1]{\left\lVert #1\right\rVert}
\newcommand{\inner}[2]{\left\langle #1,#2\right\rangle}
\newcommand{\grad}{\nabla}
\DeclareMathOperator*{\argmin}{arg\,min}
\DeclareMathOperator*{\argmax}{arg\,max}
\DeclareMathOperator{\rank}{rank}
\DeclareMathOperator{\Tr}{Tr}
\tcbset{seminarbox/.style={enhanced,breakable,sharp corners,boxrule=0pt,frame hidden,left=3.2mm,right=3mm,top=1.8mm,bottom=1.8mm,before={\par\Needspace{5\baselineskip}},before skip=0.8em,after skip=0.8em,fonttitle=\bfseries\small,coltitle=black,attach title to upper={\quad}}}
\newtcolorbox{keybox}[1][]{seminarbox,colback=NoteBlueLight,borderline west={1.8pt}{0pt}{NoteBlue},title={Key point#1}}
\newtcolorbox{paperbox}[1][]{seminarbox,colback=NoteGrayLight,borderline west={1.8pt}{0pt}{NoteGray},title={From the original paper#1}}
\newtcolorbox{suppbox}[1][]{seminarbox,colback=NoteWarmLight,borderline west={1.8pt}{0pt}{NoteWarm},title={Supplementary derivation#1}}
\newtcolorbox{warningbox}[1][]{seminarbox,colback=NoteGrayLight,borderline west={1.8pt}{0pt}{NoteGray},title={Caution#1}}
\newtcolorbox{presentbox}{seminarbox,colback=white,borderline west={2.2pt}{0pt}{black!70},fontupper=\footnotesize,top=1.2mm,bottom=1.2mm,before={\par},before skip=0.5em,after skip=0.5em,title={How to say it in the presentation}}
\theoremstyle{definition}
\newtheorem{definition}{Definition}[chapter]
\newtheorem{proposition}{Proposition}[chapter]
\newtheorem{theorem}{Theorem}[chapter]
\begin{document}
\begin{titlingpage}
\thispagestyle{empty}
\vspace*{1.2cm}
{\small\bfseries LLM Theory Seminar \hfill Week 6\par}
\vspace{1.4cm}
{\HUGE\bfseries In-Context Learning Theory\par}
\vspace{0.55cm}
{\large English seminar study notes\par}
\vspace{0.9cm}
\rule{\textwidth}{0.7pt}
\vspace{1.1cm}
\begin{minipage}{0.86\textwidth}
\small
These notes are an English companion to the seminar handout.  The emphasis is on
mathematical derivations that can be followed line by line, on separating theorem
statements from interpretation, and on identifying the assumptions under which each
claim is valid.
\end{minipage}
\vfill
{\small\textbf{Primary readings}\\[0.25em]
von Oswald et al., \textbf{Transformers Learn In-Context by Gradient Descent}.\\Aky\"urek et al., \textbf{What Learning Algorithm Is In-Context Learning?}\\Xie et al., \textbf{An Explanation of In-Context Learning as Implicit Bayesian Inference}.
\par}
\vspace{0.8cm}
{\small September 2026\par}
\end{titlingpage}
\tableofcontents*
\clearpage

\chapter{Problem setup and notation}
\section{What in-context learning means mathematically}
An in-context learner receives a prompt containing demonstrations
\[
D=\{(x_1,y_1),\ldots,(x_n,y_n)\}
\]
and a query $x_{n+1}$, then predicts
\[
\hat y_{n+1}=F_\Theta(D,x_{n+1})
\]
without changing the model parameters $\Theta$ at test time.  The central theoretical question is what computation the forward pass performs on $D$.

\section{Three distinct explanatory levels}
\begin{enumerate}
\item \textbf{Optimization view:} the forward pass implements steps of an algorithm such as GD or ridge regression.
\item \textbf{Bayesian view:} the prompt updates a posterior over latent tasks or parameters.
\item \textbf{Meta-learning view:} pretraining amortizes a task-adaptation procedure into the activations.
\end{enumerate}
These are not mutually exclusive; a posterior mean can coincide with a regularized optimizer, and a meta-trained network can implement that computation.

\begin{presentbox}
Do not ask ``Which interpretation is the one true mechanism?'' before fixing the data model, architecture, and training distribution.  Different theorems establish different levels of equivalence.
\end{presentbox}

\chapter{Linear regression as a common test bed}
\section{Ordinary least squares}
Let $X\in\R^{n\times d}$ have rows $x_i^\top$ and $y\in\R^n$.  OLS minimizes
\[
L(w)=\frac12\|Xw-y\|^2.
\]
If $X^\top X$ is invertible,
\[
\hat w_{\mathrm{OLS}}=(X^\top X)^{-1}X^\top y.
\]
A query prediction is $\hat y_q=x_q^\top\hat w_{\mathrm{OLS}}$.

\section{Ridge regression}
With penalty $\frac\lambda2\|w\|^2$,
\[
L_\lambda(w)=\frac12\|Xw-y\|^2+\frac\lambda2\|w\|^2,
\]
and the stationary condition gives
\[
(X^\top X+\lambda I)w=X^\top y,
\]
so
\[
\boxed{\hat w_{\mathrm{ridge}}=(X^\top X+\lambda I)^{-1}X^\top y.}
\]

\section{Gradient descent toward OLS}
A GD step is
\[
w^{(t+1)}=w^{(t)}-\eta X^\top(Xw^{(t)}-y).
\]
If $0<\eta<2/\lambda_{\max}(X^\top X)$ and the solution is identifiable, the iterates converge to the least-squares solution from appropriate initialization.

\section{Why linear regression is useful}
The task is simple enough that the reference algorithms are explicit, but rich enough to distinguish nearest-neighbor heuristics, one-step GD, multi-step GD, ridge regression, and Bayesian posterior prediction.

\begin{presentbox}
Linear regression provides a controlled environment where ``the Transformer behaves like algorithm A'' can be tested against an exact numerical target rather than by qualitative analogy.
\end{presentbox}

\chapter{von Oswald et al.: attention can implement a GD step}
\section{Reference update}
For examples $(x_i,y_i)$ and a current linear predictor $W$, one gradient step on squared error is
\[
W^+=W-\frac\eta N\sum_{i=1}^N(Wx_i-y_i)x_i^\top.
\]
On a query $x_q$,
\begin{align*}
W^+x_q
&=Wx_q-\frac\eta N\sum_i(Wx_i-y_i)(x_i^\top x_q).
\end{align*}
The update is a weighted sum of training residuals, with weights given by inner products with the query.

\section{Linear self-attention has the same algebra}
A linear-attention head computes a term of the form
\[
PVK^\top q=P\sum_i v_i(k_i^\top q).
\]
Choose keys to contain $x_i$, the query to contain $x_q$, and values to contain the residual-like quantity $Wx_i-y_i$.  With a suitable output projection $P$, the head can produce
\[
-\frac\eta N\sum_i(Wx_i-y_i)(x_i^\top x_q),
\]
which is exactly the correction term in $W^+x_q$.

\section{Data updates versus parameter updates}
The Transformer does not literally overwrite a persistent matrix $W$ in its external parameters.  Instead, the activations carry sufficient statistics so that the next-token prediction is the same as if $W$ had been updated.  This is a functional equivalence, not a claim that the model parameter tensor changes during ICL.

\begin{warningbox}
The clean exact construction uses linearized/linear attention and specially structured representations.  It is an existence result.  It does not by itself prove that a standard softmax Transformer trained by SGD discovers exactly this circuit.
\end{warningbox}

\begin{presentbox}
The key algebraic match is simple: GD on linear regression is a sum of residuals times query--example inner products, and a linear-attention head computes exactly a sum of values weighted by query--key inner products.
\end{presentbox}

\chapter{Aky\"urek et al.: which algorithm is implemented?}
\section{From one step to an implicit learner}
Aky\"urek et al. ask a stronger question than representability: when Transformers are trained on linear-regression sequences, which classical learning algorithm best matches their predictions?  Candidate algorithms include GD, ridge/least-squares estimators, and recursive estimators.

\section{Recursive least squares}
Let
\[
P_n=(X_n^\top X_n+\lambda I)^{-1}.
\]
Adding one example $x_n$ updates the inverse by Sherman--Morrison:
\[
P_n
=P_{n-1}
-\frac{P_{n-1}x_nx_n^\top P_{n-1}}
{1+x_n^\top P_{n-1}x_n}.
\]
Define
\[
k_n=\frac{P_{n-1}x_n}{1+x_n^\top P_{n-1}x_n}.
\]
Then the ridge/recursive estimate can be written
\[
\boxed{w_n=w_{n-1}+k_n(y_n-x_n^\top w_{n-1}).}
\]
This resembles an adaptive gradient step whose preconditioner depends on the accumulated data geometry.

\section{Depth and algorithmic complexity}
A deeper Transformer can compose multiple elementary update-like operations.  Thus architectural depth can correspond to the number or sophistication of inner learning steps, although the exact correspondence depends on the construction and training regime.

\begin{presentbox}
The important shift is from ``attention resembles gradient descent'' to ``the model can represent a family of classical estimators and may interpolate between them depending on training conditions.''
\end{presentbox}

\chapter{Bayesian linear regression}
\section{Prior and likelihood}
Assume
\[
w\sim\mathcal N(0,\tau^2I),
\qquad
y\mid X,w\sim\mathcal N(Xw,\sigma^2I).
\]
The log posterior, up to a constant, is
\[
-\frac1{2\sigma^2}\|y-Xw\|^2-\frac1{2\tau^2}\|w\|^2.
\]
Collect quadratic terms:
\[
-\frac12 w^\top\left(\sigma^{-2}X^\top X+\tau^{-2}I\right)w
+\sigma^{-2}w^\top X^\top y.
\]
Therefore the posterior is Gaussian with
\[
\boxed{\Sigma
=\left(\sigma^{-2}X^\top X+\tau^{-2}I\right)^{-1}}
\]
and
\[
\boxed{\mu=\Sigma\sigma^{-2}X^\top y.}
\]

\section{Connection to ridge regression}
Multiply the posterior-mean expression by $\sigma^2$ inside the inverse:
\[
\mu
=\left(X^\top X+\frac{\sigma^2}{\tau^2}I\right)^{-1}X^\top y.
\]
Thus the posterior mean equals ridge regression with
\[
\lambda=\frac{\sigma^2}{\tau^2}.
\]
Optimization and Bayesian inference can therefore make identical point predictions under this Gaussian model while assigning different interpretations to the same formula.

\section{Posterior predictive}
For a new $x_q$,
\[
y_q\mid x_q,D
\sim\mathcal N\left(x_q^\top\mu,
\sigma^2+x_q^\top\Sigma x_q\right).
\]
The mean is the ridge predictor; the variance expresses uncertainty that a pure point-estimation view does not encode.

\begin{presentbox}
In Gaussian linear regression, ``Bayesian ICL'' and ``regularized optimization'' can coincide exactly at the level of the posterior mean.  The distinction appears in the assumed generative model and in uncertainty, not necessarily in the point prediction.
\end{presentbox}

\chapter{Xie et al.: ICL as latent-concept inference}
\section{Pretraining distribution with a latent task}
Let $\theta$ index a latent concept or task.  A pretraining sequence is generated by
\[
p_{\mathrm{pre}}(o_{1:T})
=\int p(o_{1:T}\mid\theta)p(\theta)\,d\theta.
\]
Observing a prompt changes the posterior
\[
p(\theta\mid o_{1:t})
\propto p(o_{1:t}\mid\theta)p(\theta).
\]
Prediction marginalizes over the latent task:
\[
p(o_{t+1}\mid o_{1:t})
=\int p(o_{t+1}\mid o_{1:t},\theta)
 p(\theta\mid o_{1:t})\,d\theta.
\]

\section{Why demonstrations change behavior without weight updates}
The prompt supplies evidence about $\theta$.  The network can therefore adapt predictions by updating an implicit posterior in activation space while leaving its global parameters fixed.

\section{A simple posterior-concentration derivation}
For a finite latent set $\Theta$ and i.i.d. observations from the true $\theta_\star$,
\[
\frac{p(\theta\mid D_n)}{p(\theta_\star\mid D_n)}
=\frac{p(\theta)}{p(\theta_\star)}
\prod_{i=1}^n
\frac{p(o_i\mid\theta)}{p(o_i\mid\theta_\star)}.
\]
Taking logarithms and dividing by $n$,
\[
\frac1n\log\frac{p(\theta\mid D_n)}{p(\theta_\star\mid D_n)}
=\frac1n\log\frac{p(\theta)}{p(\theta_\star)}
+\frac1n\sum_i\log\frac{p(o_i\mid\theta)}{p(o_i\mid\theta_\star)}.
\]
By the law of large numbers, the second term converges to
\[
-D_{\mathrm{KL}}(p_{\theta_\star}\|p_\theta),
\]
so incorrect identifiable concepts lose posterior mass exponentially.

\begin{warningbox}
This i.i.d. derivation is pedagogical.  Xie et al. analyze structured sequence models such as latent concepts in HMM-like settings; the exact theorem assumptions are richer than the toy calculation above.
\end{warningbox}

\begin{presentbox}
The Bayesian account says that the prompt is evidence about which task generated the sequence.  ICL is then posterior adaptation in latent-task space rather than explicit parameter optimization.
\end{presentbox}

\chapter{Meta-learning unifies the views}
\section{Outer loop and inner loop}
Let $\tau$ denote a task sampled from a task distribution.  The outer loop optimizes shared parameters
\[
\Theta^*=\argmin_\Theta\E_{\tau}\bigl[\mathcal L_\tau(F_\Theta(D_\tau))\bigr].
\]
At test time, the inner adaptation is performed by the forward computation
\[
D_\tau\mapsto F_{\Theta^*}(D_\tau,\cdot),
\]
without an explicit gradient update to $\Theta^*$.

\section{Amortized inference}
If solving each task separately would require an optimization or Bayesian inference procedure $A(D)$, meta-training can learn a network that approximates
\[
F_\Theta(D)\approx A(D).
\]
The expensive per-task computation has been amortized into shared parameters and a fixed number of forward operations.

\section{Three views can be simultaneously true}
Suppose the task prior is Gaussian linear regression.  The posterior mean equals ridge regression.  A meta-trained Transformer can learn to compute that ridge estimator from demonstrations.  Then the same forward pass is simultaneously:
\begin{itemize}
\item a learned meta-algorithm;
\item an optimizer-equivalent computation;
\item a Bayesian posterior-mean computation.
\end{itemize}

\begin{presentbox}
Meta-learning is the outer explanatory layer: it tells us why a fixed network might acquire an inner algorithm.  Optimization and Bayesian theories characterize what that inner algorithm computes under particular task distributions.
\end{presentbox}

\chapter{What is theorem and what is interpretation?}
\section{Can implement versus does implement}
A constructive theorem of the form
\[
\exists\Theta:\quad F_\Theta(D,x)=A(D,x)
\]
shows representability.  It does not prove
\[
\Theta_{\mathrm{SGD}}\approx\Theta
\]
for ordinary training.  Empirical matching, mechanistic analysis, or an optimization theorem is needed for that stronger claim.

\section{Linear attention versus softmax attention}
Exact algebraic GD constructions often use linear or specially parameterized attention.  Standard LLMs use softmax attention with normalization, residual pathways, layer normalization, and many nonlinear blocks.  A construction may still illuminate their behavior, but exact equivalence should not be claimed without a theorem covering that architecture.

\section{From toy regression to natural-language ICL}
Natural-language prompts add latent semantics, discrete tokenization, ambiguous task definitions, distribution shift, and long-range dependencies.  The controlled linear-regression theory isolates mechanisms; it does not by itself establish a complete theory of instruction following.

\begin{presentbox}
The safest claim is mechanism-specific: ``this architecture can implement this learning algorithm under these assumptions,'' or ``trained predictions closely match this estimator on this task family.''  Avoid upgrading either statement into a universal claim about all LLM ICL.
\end{presentbox}

\chapter{Presentation checklist}
\section{Twelve formulas to keep}
\begin{align*}
&\hat w_{\mathrm{OLS}}=(X^\top X)^{-1}X^\top y,\\
&\hat w_{\mathrm{ridge}}=(X^\top X+\lambda I)^{-1}X^\top y,\\
&w^+=w-\eta X^\top(Xw-y),\\
&W^+x_q=Wx_q-\frac\eta N\sum_i(Wx_i-y_i)(x_i^\top x_q),\\
&PVK^\top q=P\sum_i v_i(k_i^\top q),\\
&k_n=\frac{P_{n-1}x_n}{1+x_n^\top P_{n-1}x_n},\\
&w_n=w_{n-1}+k_n(y_n-x_n^\top w_{n-1}),\\
&\Sigma=(\sigma^{-2}X^\top X+\tau^{-2}I)^{-1},\\
&\mu=\Sigma\sigma^{-2}X^\top y,\\
&\lambda=\sigma^2/\tau^2,\\
&p(\theta\mid D)\propto p(D\mid\theta)p(\theta),\\
&F_\Theta(D)\approx A(D).
\end{align*}

\section{Common mistakes}
\begin{itemize}
\item Confusing activation-space adaptation with an actual update of model weights.
\item Treating a linear-attention construction as an exact theorem for every softmax LLM.
\item Calling posterior mean and full Bayesian prediction the same object.
\item Inferring the trained mechanism from representability alone.
\item Assuming one toy task identifies a universal ICL algorithm.
\end{itemize}

\begin{presentbox}
Thirty-second conclusion: ICL can be formalized as an inner learning algorithm executed by a fixed network.  In controlled settings, attention can implement GD-like updates, trained Transformers can resemble classical estimators, and Bayesian posterior inference gives an alternative but sometimes algebraically identical description.
\end{presentbox}

\chapter*{Bibliography}
\addcontentsline{toc}{chapter}{Bibliography}
\begin{itemize}
\item Garg, S. et al. \emph{What Can Transformers Learn In-Context? A Case Study of Simple Function Classes}. NeurIPS, 2022.
\item von Oswald, J. et al. \emph{Transformers Learn In-Context by Gradient Descent}. ICML, 2023.
\item Aky\"urek, E. et al. \emph{What Learning Algorithm Is In-Context Learning? Investigations with Linear Models}. ICLR, 2023.
\item Xie, S. et al. \emph{An Explanation of In-Context Learning as Implicit Bayesian Inference}. ICLR, 2022.
\end{itemize}
\end{document}
